\documentclass{article}
\usepackage{amsmath}
\usepackage[hidelinks]{hyperref}
\title{Identify the loaded transition, without quietly removing the load}
\date{}
\begin{document}
\maketitle

\paragraph{Target and pinned effects.}
A warehouse operator must prepare one of two route reserves
(stay in the initial lane or detour) before the next loaded job.
The objective is to avoid normalized slower fallback; the original
delivery and ordinary side work may not be waived. The new
diagnostic's opportunity cost is the foregone optional report
service at a separate receiver, not the removal of the loaded job.
This materializes a premise of the earlier load-regime transport
calculation. Pinned \verb|SOURCE.tsv| entries are the actual selected
RWARE action and SubT dispatch, window and application callback
fragments. Initial carrier and obstacle layouts, experiment slots,
loss probabilities, motion planner adapter, fixed train/target
regime and payoff are stipulated. This is a civil reference fixture.

Let $X$ be an independently fair initial lane and $M$ a latent
pallet-at-entry condition, independently generated with probability
$p=1/4$ or $p=3/4$. One law stays fixed over independent calibration
and future mission episodes. The carrier attempts its actual
native forward step. While carrying its required load, a pallet
blocks the stock transition; the adapter then routes around to the
opposite lane using stock actions. Without the load the same
occupied shelf can be traversed, leaving the carrier in its
initial lane after the checked route. The actual resulting lane is
$Y=X\mathbin{\mathrm{xor}} M$ when loaded and $Y=X$ when unloaded.
The adapted choice depends on whether the physical move happened,
not on a forecast label. The simulator executes the selected
moves and verifies $Y$, continued load status, independent side
motion and a separate side callback.

\paragraph{Precisely what remains hidden.}
The loaded old packet contributes a three-byte debit at the
first receiving node. Each initial or final lane report is two
bytes in the same five-byte rolling window. An observer can
request one or try both; it gets at most one. A second rejected
delivery may still consume TX. Acquisition cuts are fixed
and the fixture pads the physical route schedule. The experimenter
sees only the listed callback/erasure transcript under the actions
it selected; it cannot read native route paths, live motion,
load-sensitive duration, packet-body debug logs, or hidden
pallet placement. Such data may identify a law in a different
architecture. Both unpaired marginals are fair for both laws.
A network decision adaptive in earlier episodes has the same
observed history distribution under both laws by induction;
attempting the second read after the first can only get the
same denial in the declared regime. The reserved-delivery
erasure used below is independent of the latent pallet condition.

Removing the load opens both observation slots, but the physically
observed unloaded pair distribution is equal in both laws and
does not target the loaded transition. It learns $Y=X$, not
whether a future loaded job is likely to detour. The loaded
counterfactual reserve preferences are opposite. Without a new
target-regime observation, the best worst-law staging
randomizes between the two preparations: either law's wrong
choice adds loss $|2p-1|=1/2$, for minimax regret $1/4$.
These are decision losses relative to the in-law preparation,
not estimated plant or safety-risk rates.

\paragraph{Obtain both pieces without changing the physical regime.}
A separate authorized receive path makes the second observation
available only if it is reserved. The optional four-byte report
occupies enough real receiving capacity to block the two-byte
diagnostic when left scheduled, even though the new sender can
transmit. Skipping that optional report supplies genuine RX
headroom; it does not release cargo or cancel the mandatory
three-byte work. The two resulting callbacks are joined within
one serialized calibration episode (episode identity is trusted
fixture state, not free distributed reconstruction). A no-delivery
case is explicit, the native sender debit is retained, and the
charged price of holding the optional slot is incurred even
when the second report is erased.

For one paid loaded pair, the full-transcript exact finite
minimax analysis enumerates all deterministic decoders. With
both reads, parity favors the matching preparation and gives
worst-law decision error $1/4$. In the separate half-erasure
control, one fair retained read lets the decoder split ties and
gives error $3/8$. Equal-prior pointwise likelihood gives the
same sharp lower bound in both cases. Multiplying by the wrong
reserve loss $1/2$, and comparing with $1/4$ without diagnostic,
the cost of one reserved calibration must be below $1/8$
without erasure or below $1/16$ with half erasure to be valuable.
At equality there is no strict gain. The real observed resource
difference is one missed optional report per allocated run;
its conversion to the normalized payoff is parameterized.
For state-dependent erasure this exact error calculation
would need a different experiment law.

\paragraph{Baseline and nonclaims.}
The finite exercise verifies 384 native plant/radio fixtures,
conditioned on the two fixed laws, six experiment modes and
erasure controls. Observational equality and the likelihood
verdict are exact counting results for these sources.
\href{https://arxiv.org/abs/1205.0858}{Controlled sensing}
already treats selecting informative experiments; its direct
same-information comparator ties here.
\href{https://proceedings.mlr.press/v286/marmarelis25a.html}
{Marmarelis et al.} bound off-policy predictions under
trajectory confounding with a sensitivity model. This is a
different intervention and visibility obstruction; no
reproduction of their general bound is claimed.
Preparations are chosen for a later episode, not after its fault
reveals the answer. No independent transfer, measured exchange
law, general R3b or field acceptance is asserted.
\end{document}
